% Generated draft: requires mathematical and source review.
\documentclass[11pt]{article}
\usepackage[margin=25mm]{geometry}
\usepackage{fontspec}
\setmainfont{Latin Modern Roman}
\usepackage{amsmath,amssymb,mathtools,unicode-math}
\setmathfont{Latin Modern Math}
\usepackage{enumitem,needspace,xurl,hyperref}
\setlength{\parindent}{0pt}
\setlength{\parskip}{5pt}
\setlength{\emergencystretch}{3em}
\newcommand{\N}{\mathbb{N}}
\newcommand{\Z}{\mathbb{Z}}
\newcommand{\Q}{\mathbb{Q}}
\newcommand{\R}{\mathbb{R}}
\newcommand{\C}{\mathbb{C}}
\newcommand{\F}{\mathbb{F}}
\newcommand{\A}{\mathbb{A}}
\newcommand{\PP}{\mathbb{P}}
\newcommand{\E}{\mathbb{E}}
\newcommand{\eps}{\varepsilon}
\newcommand{\dd}{\,\mathrm{d}}
\begin{document}
\section*{UnsolvedMath ID: 30006650}
\textbf{Source problem code:} catalogue entry
\subsection{Definitions and mathematical statement}
Let \(p\) be a prime. A finite \(p\)-group is a finite group \(G\) with \(|G|=p^n\) for some integer \(n\geq 0\). Denote by \(H^{\bullet\bullet}(G)\) the permutation-twisted cohomology ring named in the supplied source record. This notation is retained from the source; the record does not supply its construction, coefficient ring or field, twisting data, grading conventions, or multiplication.

\textbf{Conjecture 9 (as supplied).} For every prime \(p\) and every finite \(p\)-group \(G\), the permutation-twisted cohomology ring \(H^{\bullet\bullet}(G)\) is Noetherian.

The meaning of Noetherianity must be fixed from the underlying construction. For a commutative ring, Noetherianity means that every ideal is finitely generated, equivalently that every ascending chain of ideals stabilizes. The supplied record does not specify whether the conjecture uses this ordinary ring convention, a graded convention involving homogeneous ideals, or another convention appropriate to its ring. Thus the displayed conjecture preserves the supplied target but is not a self-contained definition of it.

\textbf{Supplied qualifications.} The source extract reports Noetherianity for every Dedekind \(p\)-group, meaning a finite \(p\)-group in which every subgroup is normal, and for every \(p\)-group of order at most \(p^3\). It attributes these cases to \([\mathrm{Mil25b},\ \text{Theorem }7.4,\ \text{Corollary }7.9]\); the record does not explicitly identify that reference bibliographically.

The extract also says, independently of Noetherianity, that an unspecified preceding theorem holds at the trivial subgroup of \(G\), yielding an alternative characterization of the usual cohomology ring of a \(p\)-group. It contains the further fragments that \(\operatorname{Spc}(K(G))\) is a Dirac scheme and that a comparison map is an open immersion. The category \(K(G)\), the comparison map, the preceding theorem, the characterization, and their precise hypotheses are absent. These contextual fragments cannot be reconstructed as precise additional assertions from the supplied record.
\subsection{Short English statement}
The supplied conjecture asks whether the permutation-twisted cohomology ring of every finite \(p\)-group is Noetherian. The record reports this for Dedekind \(p\)-groups and groups of order at most \(p^3\). The ring's construction and the intended Noetherianity convention are missing, so a self-contained formulation requires source review. The open status is supplied dated curation.
\subsection{Sources}
\begin{itemize}
\item Local UnsolvedMath registry, supplied record with permanent ID \(30006650\), exact source code \texttt{catalogue entry}, title \emph{Noetherian Twisted Cohomology of Finite \(p\)-Groups}. The record identifies the statement as Conjecture 9 and supplies the year 2026; no individual proposer is named.
\item \emph{Cohomology of Finite Groups: Interactions and Applications} (2026), Oberwolfach Reports, pp.~541--610, DOI \texttt{10.4171/owr/2026/9}. Supplied source report and literature-source label: originating Oberwolfach report. \url{https://doi.org/10.4171/owr/2026/9}.
\item Miller, \emph{Permutation twisted cohomology, remixed}. Supplied literature source: \url{https://arxiv.org/abs/2509.00954}.
\item Balmer--Gallauer, \emph{The tt-geometry of permutation modules}. Supplied literature source: \url{https://arxiv.org/abs/2307.04398}.
\item The extract cites \([\mathrm{Mil25b},\ \text{Theorem }7.4,\ \text{Corollary }7.9]\) for the stated special cases. Its identification with any listed literature source is not explicit in the supplied record.
\item Literature assessment dated exactly \texttt{2026-08-22}: the supplied curation labels the conjecture open and says that current work extends the construction and proves control results without establishing the required general Noetherianity. This is supplied dated curation, not independently verified literature assessment.
\item The registry labels the statement \texttt{corrected\_verified} and says it was reconstructed and checked against the official Oberwolfach report. These are supplied provenance claims; no independent source verification is asserted here.
\end{itemize}
\end{document}
